% Generated by tex/build_swebok_structure.py from tex/chapters/ch01/06_要求仕様化.tex
\subsection{要求仕様化}

\term{要求仕様化}とは、要求を記録し、記憶し、伝達できる形にする活動である。仕様化の方法は一つではない。自然言語で書く、構造化したテンプレートで書く、受け入れ基準として書く、モデルで表すなど、複数の方法がある。

どの方法を選ぶかは、プロジェクトの性質によって変わる。業務領域への理解、リスクの大きさ、チームの地理的分散、外注の有無、契約上の要求、要求に基づくテストの必要性、保守期間中の担当者交代などを考慮する必要がある。

\MiniBox{基本方針}{要求仕様は、読む人を意識して作る。利用者、顧客、設計者、開発者、テスト担当、保守担当は、それぞれ必要とする情報が違う。}
